\chapter{Preliminaries and State of the Art}
\label{Grundlagen}
%Verwandte Arbeiten und Forschungslücke

\section{Classification of EMG Signals}

\subsection{Motivation}

To enable patients with reduced mobility in the lower limbs to perform everyday activities as independently as possible, an assistive device must be intuitive and easy to control. Movements such as walking or climbing stairs should be supported without requiring the user to actively operate the device.
For the intent recognition in an orthosis, which means determining the movement the user intends to perform, there are multiple options available \cite{cimolato_emg-driven_2022}.\\

Many commercial solutions on the market rely on embedded mechanical sensors such as inertial measurement units, accelerometers, encoders, or load cells. These sensors provide information about metrics like limb position, velocity, acceleration, joint angle, or interaction forces and torque. Based on these measurements, the current state of the limb can be estimated and the intended movement can be inferred from gait patterns and motion dynamics. However, these approaches do not always work reliably enough, and users reportedly complain about limited accuracy and intuitiveness, which limits the suitability of such devices for everyday use \cite{cimolato_emg-driven_2022}.

This issue arises because mechanical sensors can only detect what the body and the device are doing and infer intended movements only based on the current limb state, but they are unable to directly measure the user's volitional muscle activation. To overcome this limitation, EMG sensors can be used to measure the electrical activity of muscles resulting from neural activation, which is directly related to the user's intended movement and originates from motor commands sent by the nervous system and the brain to control the muscles. In this way, EMG sensors promise to increase intuitive controllability and user acceptance \cite{cimolato_emg-driven_2022}.




\subsection{Theory}
\label{sec:EMG-Theory}

Surface electromyography sensors measure electrical potentials generated by the activation of muscle fibers. Since the signal is only measured non-invasively on the skin surface, individual muscle fibers cannot be isolated. Instead, the measured signal is a linear combination of all action potentials generated by multiple muscle fibers within the detection area of the measuring electrodes \cite{jaramillo-yanez_real-time_2020}.
Since EMG signals are noisy, highly variable and subject-dependent, further processing is required to transform the recorded signals into control information for an assistive device. The pipeline of EMG-based movement classification generally consists of six stages, although some stages may be optional depending on the application: data acquisition, segmentation, preprocessing, feature extraction, classification and postprocessing \cite{jaramillo-yanez_real-time_2020}.

For data acquisition, a review of EMG-based systems in \cite{jaramillo-yanez_real-time_2020} summarizes configurations using up to 128 EMG sensors with sampling rates between 200 and 2048 Hz. Multichannel configurations with 8 to 16 sensors are most common because they capture spatially distributed muscle activation patterns and are also available in commercial devices. Several studies use sampling frequencies around 1000 Hz, since 95\% of the relevant information in EMG signals is located below 500 Hz and, according to the Nyquist sampling theorem, the sampling rate must be at least twice the highest frequency to be captured \cite{jaramillo-yanez_real-time_2020}.

The first step applied to the raw signals is to divide them into segments for later classification. Each segment contains a time series of multiple successive samples. Different approaches either use the duration of a specific movement as a segment or apply adjacent or overlapping sliding windows on the signal stream. Larger windows can increase classification accuracy because they contain more information, but they also increase the delay before a classification result becomes available \cite{jaramillo-yanez_real-time_2020}.

During preprocessing, various techniques from classical signal processing can be applied to reduce sensor noise and improve signal quality. Common methods include offset compensation, pre-smoothing, rectification, amplification or filtering of unwanted frequency components, such as powerline interference \cite{jaramillo-yanez_real-time_2020}.

Feature extraction techniques transform the preprocessed EMG signal segments into input features for a classifier. One of the most commonly used features is the mean absolute value. Other possible time-domain features include the \ac{RMS}, variance, and standard deviation. Depending on the classifier, application, and properties of the signal, frequency-domain, time-frequency-domain, spatial, or fractal-domain features can also be used, either individually or in combination \cite{jaramillo-yanez_real-time_2020}.

To classify the extracted features and generate labels corresponding to gestures, motions, or movement intentions, various machine-learning and deep-learning methods can be used. Common approaches include support vector machines, linear discriminant analysis, and feed-forward neural networks \cite{jaramillo-yanez_real-time_2020}. Several promising state-of-the-art classifiers are discussed in the following Chapter~\ref{stateOfTheArt}.  

To further improve the quality and stability of the classification, postprocessing methods can be applied to the output of the classifier. Most commonly, majority voting is used to eliminate isolated misclassifications and outliers, while other approaches include the elimination of consecutive repetitions, thresholding, and velocity ramps \cite{jaramillo-yanez_real-time_2020}.


\subsection{State of the Art}
\label{stateOfTheArt}
A large number of classifiers and pattern-recognition methods are generally suitable for classifying streams of EMG signals. Since the goal of this work is to use the classifier in an embedded architecture for a wearable application, the focus is restricted to comparatively lightweight methods that offer low energy consumption and computational cost while providing high throughput. This requirement can make relatively simple machine-learning methods more suitable than computationally demanding deep-learning models, even if this comes at the cost of slightly lower classification accuracy \cite{castruita-lopez_electromyography_2025}.

Much of the literature on EMG classification focuses on hand gesture recognition rather than lower-limb movements. However, the underlying challenges of acquiring, processing, and classifying EMG signals are largely shared between both applications. As a result, general findings regarding suitable classifiers, EMG processing, and computational complexity can be transferred to the classification of foot movements and provide a solid methodological basis for this work.


\subsubsection{Machine Learning}

In recent work, after feature extraction, EMG classification is often performed by supervised machine-learning methods. These methods can achieve high classification accuracy while keeping computational requirements comparatively low. This is particularly important for wearable applications, where the model must run locally with limited memory and energy while providing low-latency predictions. Therefore, the relevant question is not only which classifier achieves the highest accuracy, but also which methods can be implemented efficiently on embedded hardware \cite{castruita-lopez_electromyography_2025}.
Among classical machine-learning approaches, \acp{SVM} and \ac{LDA} stand out due to their comparatively low computational requirements, which also explains their frequent use in the literature.
Other commonly investigated methods include k-nearest-neighbors, Bayesian classifiers and decision-tree-based approaches such as random forest or CatBoost \cite{castruita-lopez_electromyography_2025}.

The reviewed literature does not identify a single classifier that consistently outperforms all alternatives. Performance strongly depends on the selected features, preprocessing methods, subject-specific characteristics, and data quality. This also makes direct comparisons between results reported in different studies non-trivial. Similarly, robustness and reliability are not always evaluated under comparable conditions and are therefore difficult to assess across different studies \cite{tamilvanan_comprehensive_2026}.

Lee et al. \cite{lee_electromyogram-based_2022} compared four different methods for the classification of ten hand gestures. They used three EMG sensors, segmented the recorded signals into windows of 250 ms, and extracted six time-domain features per sensor, resulting in 18 features for classification. Although the exact results vary across the ten subjects, the mean accuracies show a clear overall trend. Logistic regression achieved the lowest accuracy with 53.9\%, while \ac{SVM} and random forest reached 87.4\% and 83.1\%, respectively. All three machine-learning methods were outperformed by an \ac{ANN}, which achieved a mean accuracy of 94.0\%. Additionally, the \ac{ANN} showed the lowest variation in accuracy between subjects, indicating the highest robustness to inter-subject variability among the compared models \cite{lee_electromyogram-based_2022}. These results motivate the consideration of neural-network-based and deep-learning methods in the following Section~\ref{sec:emgDL}.


\subsubsection{Deep Learning}
\label{sec:emgDL}

Since EMG signals are noisy, nonlinear, and subject-dependent, deep-learning methods are strong candidates for EMG classification. They are able to learn complex representations of the input data and can combine feature extraction and classification within a single learned model, reducing the dependence on handcrafted features. Deep neural networks can learn which characteristics of the input are relevant for the classification task instead of relying entirely on predefined feature-extraction methods such as those described in Section~\ref{sec:EMG-Theory} \cite{zhong_deep_2026}. The following overview is organized according to five groups of approaches that are most frequently investigated for EMG classification in recent literature.\\

%\subsubsubsection{ANN}
\ac{ANN}s, including fully connected and feed-forward neural networks, represent a comparatively straightforward approach for EMG classification \cite{li_gesture_2021}. Their structure usually consists of an input layer, one or more hidden layers, and an output layer. The layers progressively transform the input features until the output layer produces the final class prediction. Because of this simple structure, they are often used for real-time EMG classification, since they offer low computational cost and latency \cite{zhang_real-time_2019}. However, compared with the following methods, conventional ANNs do not explicitly model temporal or spatial dependencies and instead learn a direct mapping from an input feature vector to a class prediction \cite{li_gesture_2021}.

%\subsubsubsection{CNN}
\acp{CNN} are the most widely used deep-learning architecture for EMG classification and have been shown in several studies to outperform conventional machine-learning models in accuracy and robustness \cite{li_gesture_2021}. Their ability to automatically learn and extract spatial features makes them particularly suitable for multichannel EMG classification. By arranging EMG measurements according to the spatial placement of the electrodes, \ac{CNN}s are able to perform image-like spatial analysis of the signals. For high-density electrode arrays, the individual electrodes can directly correspond to pixels, while sparse multichannel signals require an appropriate spatial arrangement first. Hidden features are then extracted by repeatedly applying spatial convolutions. \ac{CNN} performance can be further increased by more complex techniques. Multi-stream CNNs first partition the EMG sensors into multiple segments based on electrode placement to learn features independently from different regions, while multi-scale \acp{CNN} use different scales for feature extraction, meaning that features are extracted using multiple kernel sizes to capture both fine-grained and more global patterns in the signal. Other techniques originating from image processing, such as attention mechanisms and vision transformers, have also been applied to advance CNN-based EMG classification \cite{zhong_deep_2026}.

%\subsubsection{Recurrent and temporal models}
Temporal models focus on the time-series characteristics of EMG rather than on spatial relations. EMG signals contain temporal muscle-activity dynamics, and dependencies between subsequent measurements provide important information about the intended movement \cite{zhong_deep_2026, li_gesture_2021}.
\acp{RNN} aim to capture these dynamics by processing the signal sequentially, one step at a time, and retaining information from previous time steps through a hidden state. Compared with simpler feed-forward networks, \acp{RNN} have been shown to achieve similar classification accuracy with fewer parameters and shorter training and inference times \cite{li_gesture_2021}. However, \acp{RNN} have difficulties capturing long-term dependencies because of the problem of vanishing or exploding gradients during training \cite{zhong_deep_2026}.
To overcome this limitation, \ac{LSTM} networks introduce a more complex memory structure, including forget gates. Together with the input and output gates, these gates control how information is updated, retained, and propagated through the network. This allows information to be preserved over longer sequences and improves the modeling of long-term temporal dependencies, which can improve classification accuracy compared with conventional \acp{RNN} \cite{zhong_deep_2026}.
As an alternative to recurrent architectures, \acp{TCN} are another important temporal model. Instead of processing a time series step by step, segments of time steps are processed in each step. \acp{TCN} apply one-dimensional convolutions along the temporal axis to each segment, similarly to \acp{CNN} on the spatial axis. Through their temporal receptive field, \acp{TCN} can capture both short- and long-term dependencies while offering competitive classification performance and greater training parallelism and scalability than recurrent architectures \cite{zhong_deep_2026}.

%\subsubsection{Hybrid and transformer models}
While \acp{CNN} primarily focus on spatial relations and temporal models exploit dependencies over time, hybrid architectures combine both approaches to capture temporal and spatial features in one model \cite{li_gesture_2021}. Hu et al. \cite{hu_novel_2018}, for example, used a \ac{CNN} for spatial feature extraction followed by an LSTM-based recurrent stage for temporal modeling. Their attention-based hybrid model achieved an improvement of up to 9.2 percentage points over previously reported state-of-the-art results \cite{hu_novel_2018}. Transformer-based methods have also been explored in hybrid architectures to combine attention mechanisms with spatial and temporal feature extraction \cite{zhong_deep_2026}.


%\subsubsubsection{Unsupervised}
The methods described above are supervised and require labeled data for training. While supervised learning remains dominant in EMG classification, unsupervised methods have achieved considerable progress in other domains such as computer vision and natural language processing. They promise to improve model generalization and adaptability and reduce the risk of human error during data labeling, which also makes them promising candidates for EMG classification \cite{zhong_deep_2026}. Frequently investigated methods include autoencoders, restricted Boltzmann machines, and deep belief networks. Studies have shown that autoencoders can outperform \acp{CNN}, while deep belief networks can outperform \ac{LDA} or \acp{SVM} in terms of classification accuracy. However, unsupervised methods can still perform worse in terms of robustness and computational efficiency, and they may increase classification errors when the data distribution changes. Overall, unsupervised learning can be considered a promising and emerging direction for EMG classification, but it has not yet become a dominant state-of-the-art approach \cite{li_gesture_2021}. But independently of the training paradigm, high classification accuracy alone is not sufficient to decide whether a method is suitable for real-world EMG control.

\subsection{Limitations in Real-World Applications}
While different approaches for EMG classification have been developed and accuracies of up to 99\% have been reported, there is no universally optimal choice of classifier, preprocessing, or feature set. Research shows continuous progress, but no single classifier has been shown to be generally superior across different EMG-control settings, especially outside laboratory conditions and offline evaluations \cite{yadav_recent_2023}.
While accuracy is the most obvious requirement, computational efficiency and robustness to changing data distributions are also essential for real-world applications such as a lower-limb orthosis. 
%subsection: computational efficiency
To enable real-time control through EMG, the user must perceive the resulting action as instantaneous after their movement intention. To meet this requirement, the computational cost of the whole pipeline, including preprocessing, feature extraction, classification, postprocessing, and action issuing, must be low enough so that the overall controller delay does not exceed 300 ms \cite{jaramillo-yanez_real-time_2020}. For lower-limb applications, even stricter timing requirements of 150 ms apply. Otherwise, the delay could be perceived as lag by the user \cite{cimolato_emg-driven_2022}.

%DL and complex models struggle to meet latency restrictions, compared to simple models with lower accuracy
Advanced deep-learning methods usually require greater computational effort than simpler machine-learning models. Performance evaluations in the literature are often conducted on computers or servers equipped with GPUs designed for computationally intensive workloads. Even though advances in model optimization have reduced inference latency, deploying these deep-learning models on energy-efficient embedded hardware is still a challenge \cite{zhong_deep_2026}.

%Wearable means small controller for low heat production and small battery -> second argument for lightweight models
For wearable devices, not only latency but also the size of the controller, heat generation, and required battery capacity influence practical usability. Balancing these requirements is challenging because they cannot all be optimized independently. To make deep-learning models compatible with constrained hardware, they often need to be converted and optimized for the target platform, potentially reducing classification accuracy and increasing latency \cite{shah_enabling_2025}.

Additionally, an optimal classifier in a real-world device must be robust to changes in EMG data distribution. Limb position, electrode placement, or physiological differences between users influence the quality and distribution of the EMG signals \cite{zhong_deep_2026}. For many of the discussed supervised models, this means that user-specific fine-tuning or retraining is required to maintain the same level of accuracy as reported in the studies. However, such retraining is inconvenient, time-consuming, and costly, which reduces the suitability of affected controllers for everyday use \cite{asfour_featureclassifier_2022}.

These limitations of current approaches motivate further research on methods that co-optimize accuracy and robustness while keeping computational cost as low as possible. A promising candidate for this purpose is Hyperdimensional Computing, which is introduced in the following section.


\section{Hyperdimensional Computing}
    \subsection{Theory and Mathematical Foundations}
    \label{sec:HDC_math_theory}

    %Intro and brain inspired
    Hyperdimensional Computing (HDC), sometimes also referred to as Vector Symbolic Architectures, is a family of computing methods in which information is represented and processed using high-dimensional distributed vectors \cite{thomas_theoretical_2021}. This computing paradigm is inspired by the human brain, not because it models individual neurons biophysically, but because it abstracts principles of neural information processing such as high-dimensional and distributed representations, parallelism, randomness, and robustness to noise. Exploiting these principles results in a computing model that can provide robustness through the distributed and partly redundant representation of information \cite{kanerva_hyperdimensional_2009}.
    
    %Mathematical Formulation
    Mathematically, an HDC model consists of a high-dimensional vector space $\mathcal{H}$, together with a set of operations defined on vectors in that space, called hypervectors \cite{schlegel_comparison_2022}. Depending on the HDC architecture, $\mathcal{H}$ can contain real-valued, bipolar, or binary hypervectors. In this work, the binary vector space $\mathcal{H}=\{0,1\}^D$ is considered, where $D$ denotes the hypervector dimension. The input data is mapped from its original input space into this high-dimensional vector space $\mathcal{H}$, which then serves as the computational domain of HDC. The idea behind this is that expanding low-dimensional input into $\mathcal{H}$ can expose useful structure for recall and learning \cite{thomas_theoretical_2021}.
    
    The usefulness of the high-dimensional space lies in three of its properties: near-orthogonality through high dimensionality, robustness through distributed representations, and computational simplicity through parallel low-precision operations.
    
    %Near-Orthogonality
    Near-orthogonality of randomly generated vectors naturally arises from high dimensionality. When a set of atomic hypervectors is sampled randomly from a high-dimensional vector space, the vectors are highly likely to be nearly orthogonal to each other. Such near-orthogonal vectors are called unrelated and are used to represent distinct entities \cite{kanerva_hyperdimensional_2009}. The probability that independently sampled random hypervectors deviate significantly from orthogonality decreases exponentially with increasing dimension \cite{hoeffding_probability_1963,thomas_theoretical_2021}.
    
    %Distributed redundancy -> ROBUSTNESS
    The robustness of HDC representations comes from the redundancy and distributed encoding of information in hypervectors. Information is not represented by a single vector element but is distributed across multiple dimensions at the same time. Due to this property, corruption of a limited number of vector elements does not necessarily destroy the represented information. This makes HDC naturally tolerant to noise in the hypervectors and to hardware errors \cite{schlegel_comparison_2022}. The same distributed representation can also provide robustness to noisy or varying input data, although the extent of this robustness depends on how the input data is encoded into the hyperdimensional space.
    
    %Fixed dimensional low-precision computation
    Computational simplicity follows from the low-precision vector elements and the operations on $\mathcal{H}$, which are described in Section \ref{sec:HDC_classification}. Most operations act independently on each vector element and can therefore be executed with a high degree of parallelism. They can also be realized in hardware using simple arithmetic or logical operations \cite{thomas_theoretical_2021}.
    

    These properties make HDC particularly interesting for embedded EMG classification, where robustness and computational efficiency are important requirements. The following section explains how HDC can be applied specifically to classification.
    
    \subsection{HDC for Classification}
    \label{sec:HDC_classification}
    % 1. Classification as prototype learning in hyperspace
    Unlike neural networks, HDC classifiers do not learn internal weights. Instead, they form one prototype vector for each class, also called the class vector. This class vector summarizes and contains all the information from the training samples belonging to this class. Classification then becomes a search for the class vector that is most similar to the encoded input sample \cite{kanerva_hyperdimensional_2009}. The following sections describe this process in detail.
    
    \subsubsection{Item Memory and Continuous Item Memory}
    % 2. Base vectors / item memory / level vectors
    The first step in HDC classification is to map input samples with multiple features each to hypervectors. Therefore, a set of base vectors, referred to as the \ac{IM}, is generated randomly, with each hypervector representing a symbol, feature identity, position or value \cite{kanerva_hyperdimensional_2009}. For categorical features and entities, it is sufficient to generate vectors randomly or pseudo-randomly because, in high dimensions, they are highly likely to be unrelated, as explained in Section \ref{sec:HDC_math_theory}, and therefore dissimilar in $\mathcal{H}$. This provides distinct representations for the corresponding entities in the vector space \cite{ge_classification_2020}. In EMG classification, such a categorical entity can be the index of an electrode, so an \ac{IM} is generated with one unrelated hypervector for each electrode.
    
    To map the numerical EMG values to $\mathcal{H}$, the representation must include not only their identity but also their ordering. Therefore, a \ac{CIM} is used that contains $L$ level hypervectors $\boldsymbol{\ell}_0, \boldsymbol{\ell}_1, \ldots, \boldsymbol{\ell}_{L-1}$, which are capable of representing $L$ distinct value levels. They are designed such that neighboring values are represented by similar hypervectors, while the distance between the hypervectors increases with the distance between the represented values \cite{ge_classification_2020}. The first vector $\boldsymbol{\ell}_0$ is generated randomly, and the following vectors are generated by successively flipping approximately $D/(L-1)$ previously unflipped vector bits per level transition, such that the total flip budget of $D$ bits is distributed as evenly as possible across the $L-1$ transitions.
    This means that consecutive level hypervectors have approximately equal distances, while increasingly distant levels become increasingly dissimilar and the first and last level hypervectors, $\boldsymbol{\ell}_0$ and $\boldsymbol{\ell}_{L-1}$, have the largest possible distance.
    
    \subsubsection{Core Operations}
    % 3. Core operations: binding, bundling, permutation
    After mapping values to the hyperspace $\mathcal{H}$, the key idea of HDC is to construct new representations by manipulating hypervectors with vector operations. For different purposes, three core operations, binding, bundling, and permutation exist \cite{kanerva_hyperdimensional_2009}.
    
    \subsubsubsection{Binding}
    % kommt near orthogonality von A*B=X auch aus der Highdimensionality? 
    %--> JA! weil d(A*B,A) = |B|
    Binding refers to a multiplication-like element-wise operation and, for binary vectors, it is realized as element-wise XOR. Binding the hypervectors $\mathbf{U}$ and $\mathbf{V}$ is denoted by $\mathbf{X} = \mathbf{U} \otimes \mathbf{V}$ \cite{ge_classification_2020}. For unrelated random hypervectors $\mathbf{U}$ and $\mathbf{V}$, the resulting hypervector $\mathbf{X}$ is near-orthogonal to both inputs and represents an association between them. Binding is used, for example, to encode address-value pairs, as described in more detail in Section \ref{sec:Encoding}. The operation is self-inverse, such that $(\mathbf{U} \otimes \mathbf{V}) \otimes \mathbf{V} = \mathbf{U}$ and binding two hypervectors with the same third hypervector preserves their distance \cite{kanerva_hyperdimensional_2009}.

    \subsubsubsection{Bundling}
    The purpose of bundling is to combine several vectors into a single representation and it is generally described as element-wise addition. For binary hypervectors $\mathbf{V}_1, \mathbf{V}_2, \ldots, \mathbf{V}_n$, the addition is followed by thresholding to obtain a binary hypervector again, which effectively results in element-wise majority voting as defined in Equation \ref{eq:binary-bundling}. Unlike binding, bundling produces a hypervector $\mathbf{Z}$ that remains similar to the input hypervectors $\mathbf{V}_i$ \cite{kanerva_hyperdimensional_2009}. Therefore, it conceptually realizes a superposition of the inputs, where several pieces of information are combined without increasing the vector dimension. Section \ref{sec:Encoding} shows how bundling is used to aggregate equally weighted hypervectors, such as feature contributions or encoded training samples, into a single hypervector. 

        \begin{equation}
            \mathbf{Z}
            =
            \bigoplus_{i=1}^{n} \mathbf{V}_i,
            \qquad
            Z_j =
            \begin{cases}
                1, & \displaystyle \sum_{i=1}^{n} V_{i,j} \geq \frac{n}{2},\\
                0, & \text{otherwise},
            \end{cases}
            \qquad
            j = 0,\ldots,D-1.
            \label{eq:binary-bundling}
        \end{equation}

    \subsubsubsection{Permutation}
    For meaningful encoding of sequences, the permutation operator is used to create a position- or order-specific version of a hypervector. In the case of EMG classification, such sequences are time series of measurements. Permutation can be expressed mathematically by multiplication with a permutation matrix and can be implemented simply as a cyclic shift of the hypervector components by one position, as defined in Equation \ref{eq:permutation}. Permutation is invertible, distance preserving, and distributive over binding and bundling \cite{kanerva_hyperdimensional_2009}. The following section shows how the three basic operations of HDC can be used to meaningfully encode samples into hypervectors. 
    
        \begin{equation}
            (\Pi \mathbf{Z})_j
            =
            Z_{(j+1)\bmod D},
            \qquad
            j = 0,\ldots,D-1.
            \label{eq:permutation}
        \end{equation}
        
    \subsubsection{Encoding}
    \label{sec:Encoding}
    % 4. Encoding: from input features to sample hypervector
    Encoding describes the process of computing a hypervector that represents a sample or a sequence of samples, including their features and corresponding values, thereby mapping the input from the feature space to the hyperdimensional space $\mathcal{H}$. Since both spatial and temporal dependencies between measurements are relevant in EMG classification, spatial encoding is first applied, followed by temporal encoding, analogously to the hybrid models described in Chapter \ref{sec:emgDL} \cite{rahimi_hyperdimensional_2016}.

    Spatial encoding of one time step can be achieved by taking, for each feature, the corresponding hypervector from the \ac{IM} and binding it to the level hypervector from the \ac{CIM} representing its value. This results in one hypervector for each feature-value pair. To combine all features of a sample into a single hypervector containing the information from the whole sample, all feature-value hypervectors of the sample are bundled together. Let $\mathbf{I}_i$ denote the IM hypervector representing feature $i$, and let $\boldsymbol{\ell}_{q_{t,i}}$ denote the CIM level hypervector representing the quantized value $q_{t,i}$ of feature $i$ at time step $t$. Then the spatial hypervector $\mathbf{Z}_t$ of one sample is given by Equation~\ref{eq:spatial_encoding} \cite{thomas_theoretical_2021,rahimi_hyperdimensional_2016}.
        
        \begin{equation}
            \mathbf{Z}_t
            =
            \bigoplus_{i=0}^{F-1}
            \left(
                \mathbf{I}_i
                \otimes
                \boldsymbol{\ell}_{q_{t,i}}
            \right).
            \label{eq:spatial_encoding}
        \end{equation}

    Temporal encoding combines $N$ consecutive spatial hypervectors $\mathbf{Z}_t$ into a single hypervector representing their temporal sequence, called an N-gram hypervector. The order of samples is encoded by applying a different number of permutations to each spatial hypervector before binding them together. Therefore, earlier samples are permuted more often than later samples. The N-gram hypervector $\mathbf{S}_t$ ending at time step $t$ is calculated as shown in Equation~\ref{eq:temporal_encoding} \cite{rahimi_hyperdimensional_2016}.

        \begin{equation}
        \mathbf{S}_t
        =
        \Pi^{N-1}\!\left(\mathbf{Z}_{t-N+1}\right)
        \otimes
        \Pi^{N-2}\!\left(\mathbf{Z}_{t-N+2}\right)
        \otimes
        \cdots
        \otimes
        \Pi\!\left(\mathbf{Z}_{t-1}\right)
        \otimes
        \mathbf{Z}_t.
        \label{eq:temporal_encoding}
        \end{equation}
    
    \subsubsection{Training}
    % 5. Training: class hypervectors / associative memory
    The training principle in HDC is class-wise accumulation. All N-gram hypervectors belonging to one class are bundled together into one hypervector, called a class vector. Due to the properties of bundling, the class vector remains similar to the N-gram hypervectors belonging to the class and therefore serves as a representative prototype for later classification. The construction of a class vector follows Equation \ref{eq:class_vector} \cite{thomas_theoretical_2021}.
    
        \begin{equation}
            \mathbf{C}_c
            =
            \bigoplus_{i=1}^{n_c}
            \mathbf{S}_{c,i}.
            \label{eq:class_vector}
        \end{equation}
    
    Here, $\mathbf{C}_c$ denotes the class vector of class $c$, $n_c$ the number of training N-grams belonging to that class, and $\mathbf{S}_{c,i}$ the $i$-th N-gram hypervector of class $c$.
    
    After processing all training samples of a class, the resulting class vector is stored in the \ac{AM} for later inference. Once this is completed for all $C$ classes, the \ac{AM} contains $C$ class vectors, which form the learned representation of the HDC classifier. However, the same encoding procedure as well as the \ac{IM} and \ac{CIM} are additionally required to encode new input samples consistently for inference \cite{ge_classification_2020}.

    
    \subsubsection{Inference}
    % 6. Inference: similarity search
    Inference in an HDC model starts by encoding the inference samples in the same way as during training. For each encoded input N-gram, the resulting N-gram hypervector $\mathbf{S}_t$ serves as a query hypervector $\mathbf{Q}$. Classification of the query hypervector is performed by comparing it with all class vectors $\mathbf{C}_c$ stored in the AM and assigning the class with the smallest distance to the query as the predicted label.
    
    As a distance measure for binary hypervectors, the Hamming distance is commonly used. It is defined as the number of vector elements that differ between two vectors, as shown in Equation \ref{eq:hamming_distance}, and can therefore be computed using simple binary operations \cite{ge_classification_2020}. The indicator function $\mathbb{1}$ evaluates to $1$ if the corresponding elements differ and to $0$ otherwise, as defined in Equation \ref{eq:hamming_indicator}.
    
    \begin{equation}
        d_H\!\left(\mathbf{Q}, \mathbf{C}_c\right)
        =
        \sum_{j=0}^{D-1}
        \mathbb{1}\!\left(Q_j \neq C_{c,j}\right).
        \label{eq:hamming_distance}
    \end{equation}

        with

        \begin{equation}
        \mathbb{1}\!\left(Q_j \neq C_{c,j}\right)
        =
        \begin{cases}
            1, & \text{if } Q_j \neq C_{c,j},\\
            0, & \text{if } Q_j = C_{c,j}.
        \end{cases}
        \label{eq:hamming_indicator}
        \end{equation}

        With the Hamming distance defined, inference consists of finding the closest of the $C$ class vectors stored in the \ac{AM}. The index of this class is assigned as the predicted label $\hat{y}$, as shown in Equation \ref{eq:inference}.

    \begin{equation}
        \hat{y}
        =
        \underset{c \in \{0,\ldots,C-1\}}{\arg\min}
        \;
        d_H\!\left(\mathbf{Q}, \mathbf{C}_c\right).
        \label{eq:inference}
    \end{equation}

        The operations and methods described above define the basic HDC classification pipeline. Their specific realization, particularly the representation and encoding of the input data, can be adapted to the requirements of different applications. The following section presents recent HDC approaches that apply and extend this basic framework for different classification tasks.

    \subsection{Related HDC Approaches}


        %\subsubsection{Image classification and general benchmark tasks}
        % MNIST / Fashion-MNIST / general HDC benchmark performance
        A common benchmark dataset for machine and deep learning in image classification is MNIST. It contains $28 \times 28$ pixel images of handwritten digits from 0 to 9 \cite{lecun_gradient-based_1998}. Hassan et al. \cite{hassan_hyper-dimensional_2022} use the MNIST dataset to compare HDC with a \ac{CNN} in terms of both accuracy and hardware complexity. The CNN implementation is based on the LeNet-5 architecture, while the HDC classifier uses spatial vector encoding. For the $28 \times 28 = 784$ pixels, the authors generate an \ac{IM} containing one random hypervector with a dimension of 10,000 for each pixel. Instead of using a second \ac{IM} to represent the pixel intensity, the corresponding pixel hypervector is permuted according to its intensity value. All resulting pixel representations are then bundled to obtain a single hypervector for each image. During training, the hypervectors of all images belonging to the same class are bundled to form the respective class vector \cite{hassan_hyper-dimensional_2022}. 
        The results show that the CNN achieves an accuracy of 99.04\%, compared with 86\% for HDC, demonstrating a large accuracy advantage of the CNN for the investigated classification task. In contrast, HDC provides clear benefits in computational efficiency. On a Raspberry Pi, HDC achieves 2.52 times faster inference and requires 2.5 times less energy than the CNN. The authors attribute the lower accuracy of HDC to the limited ability of the basic encoding scheme to extract detailed spatial features compared with CNNs \cite{hassan_hyperdimensional_2023}. 
        More advanced HDC image-encoding methods suggest that this accuracy limitation is not inherent to HDC. Smets et al. \cite{smets_encoding_2024} use point-of-interest selection and patch encoding to better capture patterns at neighboring spatial locations while relying on standard HDC operations only. Their method reaches an accuracy of 97.92\% on MNIST, substantially reducing the performance gap between HDC and CNN-based approaches \cite{smets_encoding_2024}.\\
        
        %\subsubsection{Sequential and time-series classification}
        % VoiceHD / speech recognition / sequence-like inputs
        A prominent example of HDC classification of one-dimensional data is VoiceHD by Imani et al. \cite{imani_voicehd_2017}. The authors use HDC for efficient spoken-letter classification on the ISOLET dataset. The voice signal is preprocessed and mapped from the frequency domain into hyperspace by dividing the frequency spectrum into 617 frequency bins that serve as features, with one corresponding hypervector in the IM for each frequency bin. The amplitude of each frequency bin is encoded by level hypervectors from a \ac{CIM}. Following the standard encoding procedure, feature-amplitude pairs are bound, all pairs of a sample are bundled to represent the whole sample, and samples belonging to one letter are accumulated and stored in the \ac{AM}.
        The initial HDC model achieves an accuracy of 88.4\%, which is increased to 93.8\% by retraining with misclassified samples. This is slightly higher than the 93.6\% accuracy of the neural network used as baseline. During inference, VoiceHD is 5.3 times faster and 11.9 times more energy-efficient than the neural network, demonstrating the efficiency advantage of HDC for this classification task \cite{imani_voicehd_2017}.\\
       
        %\subsubsection{HDC for EMG gesture recognition}
        While the previously discussed examples show the general utility of HDC but do not explicitly encode temporal dependencies between consecutive samples, Rahimi et al. \cite{rahimi_hyperdimensional_2016} apply HDC to EMG-based hand gesture recognition using temporal encoding. Their work forms an important foundation for the foot-movement classification model used in this thesis. 
        The authors use an EMG dataset with five classes, consisting of four hand gestures and a rest state, measured by four electrodes. The signal amplitude is discretized into 21 levels and represented by a \ac{CIM}. The encoding follows the spatial and temporal encoding principles described in Section \ref{sec:Encoding}. Channel identities are bound with their corresponding signal levels, the resulting feature-value pairs are bundled into one spatial hypervector, and consecutive samples are combined by permutation and binding to form an N-gram.
        With this HDC model, an average accuracy of  97.8\% is achieved, outperforming the SVM baseline by 8.1 percentage points. The study further shows that HDC needs less than one third of the training data required by the \ac{SVM} to reach comparable accuracy. Additionally, experiments comparing spatial encoding only with combined spatial and temporal encoding show an accuracy advantage of 7 percentage points for the spatiotemporal encoding \cite{rahimi_hyperdimensional_2016}.\\
        
        %\subsubsection{HDC for wearable high-density EMG systems}
        % Moin et al. 64-channel flexible sensor system
        Building on this work in a more realistic experimental setting, Moin et al. \cite{moin_emg_2018} evaluate HDC as part of a wearable high-density EMG acquisition system with the aim of achieving high accuracy and robustness to sensor variability. The authors use a large-area, high-density sensor array with 64 electrodes arranged in a $16 \times 4$ grid together with a robust HDC classifier. The classifier uses spatiotemporal encoding, with five consecutive samples forming one N-gram. Compared with Rahimi et al. \cite{rahimi_hyperdimensional_2016}, this setup is designed to reduce dependence on precise electrode placement by ensuring broad and stable muscle coverage. To assess this robustness, the dataset includes recordings from three different days per subject, with the electrode array being reapplied between sessions. 
        The system achieves an average accuracy of 96.64\% when trained and evaluated within the same session. Across sessions, the accuracy of HDC degrades by 7 percentage points. In comparison, the authors refer to previous SVM-based EMG studies reporting accuracy decreases of more than 30\% across sessions. These results show that the combination of high-density EMG sensors and HDC can improve robustness against session-dependent signal and electrode variability, supporting its suitability for wearable EMG control \cite{moin_emg_2018}.\\

        
        % transition:
        The discussed applications show that HDC can achieve competitive accuracy in several domains, particularly in biosignal and time-series classification. At the same time, the results also indicate that the main advantage of HDC is not necessarily superior accuracy in every task, but the combination of acceptable classification performance with simple high-dimensional operations. Since this work aims to develop a hardware-efficient classifier, the following section discusses why HDC is particularly attractive for embedded and edge devices.

    
    \subsection{HDC for Embedded and Edge Devices}
    \label{sec:HDC_for_hardware}
    The related approaches discussed above show that HDC is not always the most accurate model, especially in image classification, but that it can provide substantial advantages in speed and energy efficiency. This raises the question of why HDC maps especially well to embedded and edge hardware.
    
    %Hardware advantages of HDC
    %1. simple operations
    The first reason is the low-precision representation of hypervectors and the simplicity of HDC operations. As already described, hypervectors consist of low-precision vector elements. Different HDC models use different element types, including binary $\{0,1\}$, bipolar $\{-1,1\}$, ternary $\{-1,0,1\}$, integer, or real-valued representations. Non-binary models can achieve higher accuracy, but binary hypervectors are especially attractive for hardware because they reduce memory requirements and allow the core operations to be implemented with parallel bit-level operations \cite{thomas_theoretical_2021,ge_classification_2020,martino_configurable_2026}. 
    For binary hypervectors, binding can be realized as element-wise XOR. Bundling requires element-wise addition or accumulation followed by a comparison against a threshold. Permutation can be implemented as a cyclic shift of vector elements, and the Hamming distance can be computed by applying the XOR operation to two hypervectors and counting the differing bits \cite{hassan_hyper-dimensional_2022}. Therefore, the basic operations of binary HDC avoid complex arithmetic such as floating-point operations, multiplication, or division, which makes them suitable for low-cost hardware.
    
    %2. high parallelism + easy algorithmic layout
    The second reason is the high degree of data-level parallelism. In a fully parallel implementation, HDC operations can in principle be applied to all vector components simultaneously, because there are no data dependencies between different dimensions. In practical architectures, the latency depends on the available vector width, memory bandwidth, and degree of hardware parallelism \cite{martino_configurable_2026}.
    This configurability exposes a direct trade-off between execution time and resource usage and is one of the main reasons why HDC is attractive for hardware optimization.

    %Short disadvantages 
    However, the same properties also introduce challenges. Large hypervectors are commonly used because they provide redundancy and improve the probability of near-orthogonality. At the same time, all vector operations, memory accesses, and communication scale with the vector dimension. Consequently, unnecessarily large hypervectors increase model size, inference time, energy consumption, and hardware resource usage, which can reduce the practical benefits of HDC on embedded devices \cite{basaklar_hypervector_2021}.\\
    
    %Optimization directions.

   This motivates optimization of both the hardware architecture and the HDC representation. On the hardware side, HDC is well suited for architectures that can directly exploit its wide vector operations. Dedicated ASICs, \acp{FPGA}, and custom processor extensions can implement the operations with extensive data parallelism and local memory structures, in contrast to scalar processor implementations that can only make limited use of this parallelism \cite{martino_configurable_2026}. 
    On the representation side, decreasing the vector dimension directly lowers the number of hypervector elements that have to be stored, transferred, and processed. Similarly, reducing the number of quantization levels decreases the size of the \ac{CIM}. Both approaches are attractive for embedded hardware, but reducing the vector dimension can decrease representational redundancy, while reducing the number of quantization levels lowers the resolution of the continuous input representation. Both effects can reduce classification accuracy.

    For this reason, this work does not simply reduce $D$ and $L$, but also aims to optimize the representation itself. The quantization strategy is adapted to map EMG values more effectively to discrete levels, and the arrangement of level hypervectors in the \ac{CIM} is designed to make the encoding more discriminative even with fewer dimensions or fewer levels. The foundations of these optimization directions are discussed in the following section.

    \subsection{Optimization of HDC Representation}
    \label{sec:optimizationOfHDCRelatedWork}
    As a foundation for the optimization methods used in this work, this chapter presents two related works that improve the performance of HDC by refining the underlying hypervector representations. 
    
    Basaklar et al. \cite{basaklar_hypervector_2021} address the issue that large vector dimensions are required to achieve high accuracy, but at the same time increase model size, inference time, energy consumption, and hardware cost. Therefore, the authors develop an approach to reduce the vector dimension by optimizing the representation of numerical feature values in the \ac{CIM}. This way, a dimension reduction by more than a factor of 32 is achieved while maintaining classification accuracy and increasing the robustness of the classifier \cite{basaklar_hypervector_2021}.
    Their method is based on changing the distances between hypervectors in the \ac{CIM}. While a conventional \ac{CIM} uses equally spaced vectors, Basaklar et al. define the number of bits flipped for each level transition as the optimization parameter, such that not all value intervals are represented with the same distance in the hyperspace $\mathcal{H}$. Through such a non-equally spaced \ac{CIM}, more important regions of the value space can be mapped more meaningfully to the hyperspace, increasing the distance between class vectors and thereby potentially improving classification accuracy and robustness. The selection of the number of bit flips for each level transition and feature is formulated as a non-convex, discrete multi-objective optimization problem, with classification accuracy and robustness as the two objectives \cite{basaklar_hypervector_2021}.
    
    To solve this optimization problem, a genetic algorithm is used. A population of different candidate solutions, here bit-flip distributions over the level transitions, is generated and iteratively modified through mutation and crossover. After evaluating the resulting individuals, the best solutions are selected to form the population of the next generation. This process is repeated for a predefined number of generations. In the HDC case, each individual contains one bit-flip vector per feature, where each vector element specifies the number of bits to flip for the corresponding level transition. Through mutation and crossover, the bit flips are redistributed across the level transitions. To evaluate an individual, the corresponding CIM is generated using the exact bit flip counts and the CIM is then used to train and evaluate an HDC model. The resulting models are evaluated with respect to both classification accuracy and class-vector similarity, which are used to select candidate solutions for the following generation \cite{basaklar_hypervector_2021}.
    
    In addition to accuracy, robustness is considered as a second optimization objective. It is measured based on the similarity between the class vectors. The more dissimilar the class vectors are, the better separated the classes are in hyperspace and the more robust the model is considered. Since the individuals are ranked according to both accuracy and robustness, there is not necessarily a single best solution. One solution may achieve higher accuracy, while another provides higher robustness. Therefore, Pareto fronts are used to rank the solutions. A solution dominates another solution if it performs better in at least one objective while being at least equally good in all other objectives. Consequently, two solutions for which one provides higher accuracy but the other higher robustness can belong to the same Pareto front. After this ranking, solutions from the best Pareto fronts are selected to form the next generation. This multi-objective optimization of the CIM forms the foundation for the optimization method presented in Chapter \ref{sec:GA} \cite{basaklar_hypervector_2021}. \\

    Another essential work in the context of this thesis is Ledwig \cite{ledwig_vdl-implementierung_2025}. He implements an HDC model in VHDL for an FPGA using the same foot-EMG classification task as this thesis and therefore provides a direct predecessor to the optimization investigated in this work. The implementation shows that HDC can be efficiently executed on an FPGA when implemented in a resource-aware manner. Apart from studying latency, efficiency, and classification accuracy depending on the N-gram length and training set size, Ledwig also analyzes the influence of the number of discretization levels $L$ on classification accuracy. The experiments reveal a non-monotonic dependence on the number of CIM levels. The highest mean accuracy is achieved with 40 levels. Increasing the number to 50 initially reduces the accuracy, while 100 levels again achieve a comparable result. Further increasing the number of levels does not provide a consistent improvement \cite{ledwig_vdl-implementierung_2025}.
    
    This behavior is attributed to the quantization of continuous feature values into discrete level hypervectors. With uniform quantization, changing $L$ also changes the positions of the level boundaries in the feature space. As a result, the classification performance depends not only on the available number of levels, but also on how well the resulting boundaries match the distribution of the input data and separate the classes in the quantized feature space. For the investigated datasets, the boundary positions obtained with 40 and 100 levels provide particularly suitable representations, while other level counts can result in less favorable mappings \cite{ledwig_vdl-implementierung_2025}. These findings motivate the investigation of whether deliberately selecting the level boundaries can provide a more meaningful mapping of the input data and thereby maintain or improve classification accuracy with fewer levels. This question forms the foundation for the quantization optimization presented in Chapter \ref{sec:quantization}.
    
   % \subsection{Limitations of HDC}
   % Before describing the methods of this work to improve the HDC model for FPGA deployment %without sacrifyzing accuracy, this section presents a theorethical analysis of what %accuracies are even reachable depending on the data and problem. Readers who are not %interested in those sidefacts can skip this and proceed with chapter \ref{sec:background}.
   % %Mathematical foundations, what is expressable, what not? 
   % %maybe move this to background
    
    %\cite{thomas_theoretical_2021}
   % \cite{yu_understanding_2022}

    \section{FPGA and High-Level Synthesis}
    \label{sec:FPGA_background}
    Apart from optimizing the HDC model for EMG classification, this work investigates its implementation as a hardware accelerator for embedded execution. Compared to running only software on a general-purpose processor, application-specific hardware can improve processing performance and energy efficiency. At the same time, hardware design is more complex because low-level behavior has to be specified in detail \cite{nane_survey_2016}. To avoid a fully manual implementation at \ac{RTL}, the accelerator is described in SystemC and synthesized to \ac{RTL} using \ac{HLS} before deployment on an \ac{FPGA}.

    %1. FPGA architecture and resources
        An \ac{FPGA} is not a sequential processor, but a configurable hardware platform mainly composed of programmable logic, routing resources, and on-chip memory. Combinational logic is primarily implemented through \acp{LUT}, which act as configurable truth tables and can realize Boolean functions such as XOR operations, comparisons, and parts of arithmetic circuits.
        Configurable interconnects connect the logic and memory resources to form larger application-specific datapaths. Registers store intermediate values and state between clock cycles, while \acp{BRAM} provide larger on-chip memories. Depending on the \ac{FPGA} model, additional components such as \ac{DSP} blocks, UltraRAM, High-Bandwidth Memory (HBM), and specialized I/O interfaces can provide additional capabilities for specific applications \cite{fang_-memory_2020}.
        
    %2. Parallelism, pipelining and dataflow
        Besides configurability for application-specific circuits, a main advantage of \ac{FPGA}s is their ability to enable a high degree of parallelism. Independent operations can be implemented as separate hardware structures and can therefore execute in the same clock cycle. Pipelining can also be realized on an FPGA by dividing a computation into stages, so that different inputs can be processed in different stages at the same time. This can substantially improve throughput compared with processing inputs sequentially \cite{fang_-memory_2020}. 

        At the same time, increasing parallelism requires more hardware resources such as \ac{LUT}s and routing resources, which are limited on a real device \cite{fang_-memory_2020}. Pipelining additionally requires registers between pipeline stages, but can shorten the combinational paths between them and thereby enable higher clock frequencies. The maximum achievable frequency is limited by the critical path, which is the longest timing path between sequential elements in the hardware design. If all timing paths satisfy the constraints defined by the target clock period, the design achieves timing closure. This is important because the operating frequency directly affects latency and throughput \cite{zheng_fast_2014}. Therefore, the achievable degrees of parallelism and pipelining are limited in practice by the available FPGA resources and the required target frequency.\\
        
    %3. SystemC and HLS with FPGA implementation flow
         Since describing hardware directly at \ac{RTL} becomes increasingly complex for large systems, \ac{HLS} can be used to avoid a fully manual low-level implementation. \ac{HLS} tools can start from synthesizable, hardware-oriented descriptions in high-level languages such as SystemC and automatically generate an \ac{RTL} description that realizes the described behavior \cite{nane_survey_2016}. However, this does not mean that SystemC code is executed as software on the \ac{FPGA}. Instead, the SystemC description serves as the behavioral input from which the actual hardware structure is derived.

        The main tasks performed during \ac{HLS} are allocation, scheduling, and assignment. Allocation determines the number and type of hardware resources used to implement the desired behavior. Scheduling decides in which clock cycle each operation is executed, and assignment binds variables and operations to concrete hardware resources such as registers and functional units \cite{cheng_novel_2007}. 
        
        The generated \ac{RTL} description can then be simulated to verify that the hardware behavior still matches the intended algorithm. After this functional check, the design is passed to the \ac{FPGA} implementation flow. During FPGA synthesis, the \ac{RTL} description is transformed and mapped to a netlist consisting of resources available on the target \ac{FPGA}, such as \acp{LUT}, registers, memories, and dedicated arithmetic blocks \cite{amd_vivado_2026}. Afterwards, placement assigns these resources to physical locations on the chip, while routing connects them through the programmable interconnect fabric \cite{amd_vivado_nodate}.
    

    %4. Hardware evaluation metrics
        The resulting hardware implementation can be evaluated using both resource and performance metrics. Resource utilization describes how much of the target \ac{FPGA} is occupied by the design. The most relevant resources are \ac{LUT}s for combinational logic, flip-flops for registers, \acp{BRAM} for larger on-chip memories, and \ac{DSP} blocks for dedicated arithmetic operations \cite{nane_survey_2016}. These metrics indicate whether the design fits on the selected device and which resource type limits further parallelization \cite{imani_fach_2019}.

        Performance is characterized by clock frequency, latency, and throughput. The clock frequency determines how fast the synchronous hardware design can be operated. The latency in clock cycles, meaning the number of cycles required until a result is produced, must also be considered, since a design with a high clock frequency can still be slow if many clock cycles are required per input sample. Throughput describes how frequently new inputs can be processed or new results can be produced \cite{nane_survey_2016}. 

    %5. Transition and relevance for HDC
        These concepts are particularly relevant for \ac{HDC}, because the basic operations on binary hypervectors are simple, expose a high degree of data-level parallelism, and can be decomposed into many independent bit- or word-level operations. This makes \ac{HDC} well suited for FPGA acceleration, where such independent operations can be mapped to parallel hardware structures. At the same time, the large dimensionality of hypervectors creates a direct trade-off between latency and resource usage. A fully parallel implementation can reduce the number of required cycles, but it also increases the amount of logic, memory bandwidth, and routing resources required. Therefore, the practical accelerator design must choose a suitable degree of parallelism for the available FPGA resources and timing constraints.

        With these hardware concepts established, Chapter~\ref{sec:background} first introduces the concrete EMG dataset and HDC baseline model used in this work. The actual SystemC accelerator architecture is then presented later in Section \ref{sec:fpga_accelerator}.





   

\chapter{Background}
\label{sec:background}
%--> what i specifically used in my project
With these hardware concepts established and before presenting the optimization methods, this chapter defines the EMG dataset and the HDC baseline model used throughout this work.

    \section{EMG Data}
    \label{sec:emg-preprocessing}
    % dataset, windows, features, labels, transitions, preprocessing etc.

    This work uses four subject-specific EMG datasets that were recorded for the development of an orthosis system for lower-limb movement assistance. The participants performed foot flexion, extension, inversion, and eversion, while the rest position was used as the fifth class. The recordings were acquired with a bracelet containing 32 electrodes (HD20MM1602B, OT Bioelettronica S.r.l., Turin, Italy), arranged in two columns and 16 rows. The bracelet was worn on the upper part of the shin below the knee, and the signals were sampled at 2000 Hz.
    After recording, the data were preprocessed using a band-pass filter between 10 and 500 Hz and a notch filter at 50 Hz to reduce noise and power-line interference. Afterwards, the \ac{RMS} was computed over overlapping windows of 504 samples with a shift of 18 samples. This preprocessing was applied to each of the 32 channels, resulting in 5,210 feature vectors with 32 RMS features per subject \cite{ledwig_vdl-implementierung_2025, cnejevici_closed-loop_2026}. The preprocessed data are provided as CSV files, with the \ac{RMS} features scaled to the range $[-1,1]$. Each dataset is divided into 3,040 training samples, 1,300 validation samples, and 870 test samples.
    
    \section{HDC Pipeline}
    %dataflow in training and classification
    The HDC baseline model that is optimized in this thesis was developed prior to this work, based on the MATLAB HDC model of Rahimi et al. for EMG-based hand gesture classification \cite{rahimi_hyperdimensional_2016}. It additionally incorporates concepts from Ledwig's HDC implementation for the same datasets \cite{ledwig_vdl-implementierung_2025}. The baseline model as well as all contributions developed in this thesis are available at \url{https://gitlab.cs.fau.de/ew14ozom/hdc-emg-optimization}. It is implemented in C and organized into multiple modules. The following sections describe these modules and their role in the HDC pipeline.
        \subsection{Data Loading and Quantization}

        The model operates on EMG samples that were preprocessed as described in Section \ref{sec:emg-preprocessing}. Each sample is represented as a feature vector with 32 RMS values, one for each electrode, and annotated with one of five class labels from 0 to 4. The quantizer uses uniform boundaries to map each feature value to a discrete level for subsequent mapping to the hyperspace $\mathcal{H}$. The number of levels $L$ is configurable as a model parameter. The quantization boundaries in the feature space are precomputed and stored in the quantizer, so that each feature value can be quantized by comparison with the stored boundaries. Although such edge cases are not present in the datasets used in this work, values below the first or above the last boundary are mapped to the first and last level, respectively. The quantization step forms the interface between the continuous EMG features and the discrete level representation used by the HDC model.

        
        \subsection{Item Memory}
        \label{sec:baseline_item_memory}
        The item memory provides the base hypervectors used to map EMG samples to the hyperspace. The module supports two modes that determine the layout of the \ac{IM} and \ac{CIM} and consequently also affect the encoding process.

        The first mode uses separate \ac{IM} and \ac{CIM}. The \ac{IM} contains 32 unrelated random hypervectors representing the electrodes, while the \ac{CIM} consists of $L$  equally-distant hypervectors representing the quantized feature values. During encoding, the corresponding hypervectors from the \ac{IM} and \ac{CIM} are bound for each feature-value pair. This corresponds to the conventional HDC encoding described in Section \ref{sec:Encoding}. However, this representation uses one shared CIM for all features and therefore does not allow the level representation of each feature to be optimized independently, which is required for the optimization described in Chapter \ref{sec:GA}.

        The second mode is inspired by an alternative implementation provided by Rahimi et al. in the code accompanying their HDC model \cite{rahimi_hyperdimensional_2016} and combines feature identity and quantized value in a \ac{CCIM}. The \ac{CCIM} consists of 32 feature-specific CIMs, one for each electrode, and can be represented as a two-dimensional array of hypervectors. Consequently, the feature identity and its quantized value are already fused in the corresponding CCIM hypervector, eliminating the binding between separate IM and CIM hypervectors from the encoding process.
        
        The CCIM is initialized as follows. For each feature, the first level hypervector is initialized randomly, resulting in 32 unrelated hypervectors across the features. In addition, a separate random permutation of the hypervector indices is generated for each feature. These permutations determine the order in which vector elements are flipped when constructing the following level hypervectors. Starting from the first level, approximately $D/(L-1)$ previously unchanged bits are flipped to generate each subsequent level hypervector, such that the total flip budget of $D$ bits is distributed as evenly as possible across the $L-1$ level transitions. In this way, each feature obtains a CIM with $L$ hypervectors that preserves locality between quantization levels, because neighboring levels differ in only a fraction of vector elements, while the Hamming distance increases for more distant levels.

        
        \subsection{Encoder}
        The encoder uses the quantizer and CCIM to map the EMG samples to the hyperspace $\mathcal{H}$. For the spatial encoding of each sample, the encoder iterates over all 32 features, quantizes the corresponding feature value, and retrieves the feature-level hypervector from the CCIM. Since the CCIM already combines the feature identity and its quantized value, no additional binding operation is required. The 32 retrieved hypervectors are bundled into one spatial hypervector representing the complete EMG sample.

        Temporal encoding combines N consecutive spatial hypervectors into one N-gram hypervector using the permutation and binding operations described in Section \ref{sec:Encoding}. The resulting N-gram represents both the spatial information of the individual samples and their temporal order.
            
        \subsection{Trainer}

        The trainer module constructs five class vectors from the training data, one for each class. It processes the training set as a stream of EMG samples and uses the encoder module to generate N-grams from this stream. N-grams that contain samples from different classes are skipped, since they cannot be assigned unambiguously to a single class. All valid N-grams belonging to the same class are bundled into one class vector. Once all five class vectors are stored in the AM, training is complete and the model can be used for inference.
        
        \subsection{Evaluator}
        The evaluator encodes the test set as a stream of EMG samples using the same encoder as during training. Each query hypervector is compared with the class vectors stored in the AM by computing the Hamming distance. The implementation maps the Hamming distance to a similarity score between $-1$ and $1$, where $1$ represents identical hypervectors and $-1$ represents hypervectors that differ in every dimension. The predicted class is determined by the class vector with the highest similarity to the query hypervector.
        
        In contrast to training, N-grams containing samples from different classes are not skipped to study the classification behavior during transitions between movements as well. Predictions for such transition N-grams are tracked separately. In addition, the evaluator records correct and incorrect predictions, the confusion matrix, overall accuracy, and class-average accuracy.
        
    \section{Baseline Model Configuration}
    The described HDC baseline model exposes multiple parameters that influence classification accuracy and model complexity. This section fixes these parameters to establish a reference configuration for the subsequent optimization.
    
      First, binary hypervectors are used throughout this work because their simple logical operations are particularly suitable for hardware implementation, as discussed in Section \ref{sec:HDC_for_hardware}. The vector dimension is set to $D$ = 10,000. This dimension is a common standard in the HDC literature \cite{kanerva_hyperdimensional_2009,rahimi_hyperdimensional_2016,imani_fach_2019,hassan_hyper-dimensional_2022,yu_understanding_2022} and therefore provides a conventional reference point for the baseline model. Whether this dimension is actually required for the investigated EMG classification task is evaluated in Chapter \ref{sec:results_baseline}.
   The number of quantization levels is set to $L$ = 40, since this configuration achieved the highest mean classification accuracy in the experiments of Ledwig \cite{ledwig_vdl-implementierung_2025}. Finally, the N-gram length is set to $N$ = 3. As shown in Section \ref{sec:results_baseline}, larger N-gram sizes provide only marginal additional accuracy in the unoptimized baseline, while $N=3$ already captures temporal context with a shorter sequence length and smaller buffer requirements.

   
    Furthermore, the baseline uses the \ac{CCIM} instead of separate \ac{CIM} and \ac{IM} structures. This choice is essential for the feature-specific \ac{CIM} optimization with a genetic algorithm described in Section \ref{sec:GA}, since each feature requires its own independently adjustable sequence of level hypervectors. Compared with separate \ac{CIM} and \ac{IM} structures, however, the \ac{CCIM} increases the required memory. With $F = 32$ features, storing the \ac{CCIM} requires $F\cdot L \cdot D$ bits. For the baseline configuration, this results in $32 \cdot 40 \cdot 10{,}000~\mathrm{bits}
= 12{,}800{,}000~\mathrm{bits}
= 1{,}600{,}000~\mathrm{bytes}
= 1.6~\mathrm{MB}$. In comparison, separate IM and CIM structures require $(F+L)\cdot D$ bits, resulting in $(32 + 40) \cdot 10{,}000
= 720{,}000~\mathrm{bits}
= 90{,}000~\mathrm{bytes}
= 90~\mathrm{kB}$. The CCIM therefore requires approximately 17.8 times more memory for the baseline configuration. 
    At the same time, the \ac{CCIM} simplifies spatial encoding because the feature identity and quantized value are already combined in the stored hypervector. Encoding one sample therefore only requires 32 \ac{CCIM} lookups, reading 320{,}000 bits, or 40~kB, of hypervector data for $D = 10{,}000$. With separate \ac{IM} and \ac{CIM} structures, 32 \ac{IM} and 32 \ac{CIM} hypervectors have to be accessed per sample, corresponding to 640{,}000 bits, or 80~kB, of data without considering caching. In addition, the separate representation requires binding each of the 32 feature-level pairs, resulting in $32 \cdot 10{,}000 = 320{,}000$ bitwise XOR operations per spatially encoded sample. These binding operations are not required with the CCIM.\\

    The higher memory requirement of the CCIM is therefore an important limitation of the baseline model and directly motivates the reduction of both $D$ and $L$. At the same time, its feature-specific structure provides the basis for independently optimizing the level representation of each EMG feature.
